%%
%% SMIF paper (Single-Molecule Imputation using Fiber-seq). Built from the acmart template of proposal.tex (same
%% class options, authors and bibliography). Unwritten sections carry a %TODO with the results they will draw on
%% (reports/Nov_13_extended.md, Nov_20.md, Nov_20_figs.md).
%%
%% Commands for TeXCount
%TC:macro \cite [option:text,text]
%TC:macro \citep [option:text,text]
%TC:macro \citet [option:text,text]
%TC:envir table 0 1
%TC:envir table* 0 1
%TC:envir tabular [ignore] word
%TC:envir displaymath 0 word
%TC:envir math 0 word
%TC:envir comment 0 0
%%
\documentclass[acmtog,nonacm]{acmart}   % nonacm: no ACM reference box, no journal footer

\AtBeginDocument{%
  \providecommand\BibTeX{{%
    Bib\TeX}}}

%% No ACM boilerplate: no copyright/permission block, no "Authors' Contact Information" footnote
\setcopyright{none}
\authorsaddresses{}

\citestyle{acmauthoryear}
\usepackage{booktabs}

\begin{document}

\title{SMIF: Single-Molecule Imputation using Fiber-seq}

\author{Azaan Rehman}
\email{ara199@sfu.ca}
\correspondingauthor
\affiliation{%
  \institution{Simon Fraser University}
  \city{Burnaby}
  \state{British Columbia}
  \country{Canada}
}

\author{Mehdi Foroozandeh}
\affiliation{%
  \institution{Simon Fraser University}
  \city{Burnaby}
  \state{British Columbia}
  \country{Canada}
}

\author{Andrew B. Stergachis}
\affiliation{%
  \institution{University of Washington}
  \city{Seattle}
  \state{Washington}
  \country{USA}
}

\author{Christine Queitsch}
\affiliation{%
  \institution{University of Washington}
  \city{Seattle}
  \state{Washington}
  \country{USA}
}

\author{Maxwell W. Libbrecht}
\affiliation{%
  \institution{Simon Fraser University}
  \city{Burnaby}
  \state{British Columbia}
  \country{Canada}
}

\renewcommand{\shortauthors}{Rehman et al.}

\begin{abstract}
Bulk genomic assays such as ChIP-seq and ATAC-seq report how a chromatin property is distributed along the
genome, but only as an average over millions of DNA molecules. Fiber-seq reads chromatin accessibility,
nucleosome positions and CpG methylation on individual molecules, but no technology measures histone marks or
protein binding on many molecules at once. We ask whether these bulk properties can be imputed onto each molecule
when the only supervision is bulk. This is a weakly supervised regression problem: each training example is a bag
of up to 100 molecules, and its label is the average over a large population. We present SMIF, which maps
every Fiber-seq read, independently and with its reference sequence as context, to eight per-base tracks (ATAC,
CTCF and six histone marks). It averages these predictions in fold-change space, so that the bag prediction can be
trained against ENCODE tracks. Because no molecule-level labels exist, we also propose an evaluation protocol in
which every test has a known value under the null that the model predicts the bulk average on every molecule. In
bulk, the model predicts cell-type-specific signal beyond a cross-cell average, and in training cell types it
reaches two thirds of the correlation that the agreement between ENCODE replicates allows. Trained only on
mixtures of two cell types, it separates their molecules as well as a model trained on each cell type's own
tracks. On single molecules, although its only CTCF supervision is bulk ChIP-seq, the model assigns more CTCF to
molecules whose CTCF motif is protected than to open molecules at the same site (83--94\% of bound sites,
replicated in a second training), and not at unbound motifs. On the X chromosome, never seen in training, it
separates the active and inactive copies at 86--90\% of silenced promoters for the active marks, but not for
H3K27me3. Two models with the same bulk accuracy can differ tenfold in these single-molecule readouts, so bulk
accuracy alone cannot validate a per-molecule imputation.
\end{abstract}

\maketitle

\section{Introduction}
\label{sec:intro}
% TODO: bulk vs single-molecule chromatin; what Fiber-seq measures and what it does not; the weak-supervision
% problem (bag-average labels, instance-level predictions); why bulk accuracy alone cannot show
% single-molecule information (the null H0); contributions.

\section{Related work}
\label{sec:related}

\textbf{Single-molecule chromatin profiling.} Fiber-seq \citep{Stergachis20} marks accessible DNA with a
non-specific adenine methyltransferase and reads each chromatin fiber end to end with long reads. One read gives
the accessibility, nucleosome positions and CpG methylation of one molecule over 5--25\,kb, and fibertools
\citep{Jha24} turns the raw calls into per-molecule features. DiMeLo-seq \citep{Altemose22} uses the same principle
to map where one chosen protein binds on single molecules, one target per experiment. Single-cell CUT\&Tag
\citep{Bartosovic21} profiles histone marks in single cells, but with sparse coverage and one or a few marks per
cell. Together, these methods show that \emph{measuring} many chromatin properties on the same molecule is still
out of reach. We instead \emph{impute} them onto Fiber-seq reads.

\textbf{Epigenome imputation.} ChromImpute \citep{Ernst15}, PREDICTD \citep{Durham18} and Avocado
\citep{Schreiber20Avocado} predict missing bulk tracks from the other tracks of the same cell type, by regression
or by tensor factorization over cell types $\times$ assays $\times$ positions. CANDI \citep{Foroozandeh25} imputes
and denoises bulk assays with self-supervision and calibrated confidence. dHICA \citep{Wen24} shows that one input
assay can be enough: it predicts several histone-mark tracks from bulk accessibility (ATAC-seq or DNase-seq) and
the DNA sequence, the bulk analogue of our input. All of these output population averages, and all need bulk data
from the target cell type. SMIF keeps the targets (ENCODE assays) but changes the unit of prediction from a track
to a molecule. Its input is accessibility measured on each molecule, so a new cell type needs only its fibers.
\citet{Schreiber20Pitfall} showed that models predicting across cell types often fail to beat the average of the
training cell types. We use that average as our bulk baseline and score the part of each track that it does not
explain.

\textbf{Learning from aggregate labels.} In multiple-instance learning \citep{Dietterich97, Ilse18}, a label
belongs to a bag of instances, and in its standard form it says whether at least one instance is positive.
Learning from label proportions \citep{Quadrianto09} uses the fraction of positive instances in each bag, and can
recover instance labels only when bags differ in composition. Our problem is a regression analogue of the second
setting: a bag is the stack of fibers over a window, its label is the bulk track, and the label is a mean over a
much larger population than the bag. Bulk deconvolution methods such as CIBERSORT \citep{Newman15} also separate a
bulk signal into its parts, but into cell-type proportions and profiles, not into molecules. The mixture test in
Section~\ref{sec:eval-mixture} uses the identifiability argument of learning from label proportions: it varies the
composition of the bags so that each population's signal is determined.

Compared with these works, we connect bulk-assay imputation to single-molecule inputs, and we give an evaluation
protocol for a setting in which no molecule-level labels exist.

\section{Problem setup}
\label{sec:problem}

We set up the problem in four stages: the genome coordinates, the molecules (instances), the bulk tracks (bag
labels), and what the model must output. Table~\ref{tab:data} summarizes the objects for one window, and
Appendix~\ref{app:bio} gives the biology background.

\textbf{Coordinates and reference.} The human genome is a string of about $3\times10^9$ letters over
$\{A,C,G,T\}$, split into chromosomes. We work on windows of $L=5{,}000$ consecutive positions, given as
(\emph{chromosome, start, end}). The reference sequence of a window enters the model one-hot,
$S\in\{0,1\}^{4\times L}$. It is the same in every cell type and every molecule, so it tells the model \emph{where}
it is, but not \emph{which} cell type, or which molecule, it is looking at.

\textbf{Molecules: the instances.} A Fiber-seq read is one chromatin fiber of a cell population. After standard
processing \citep{Jha24}, each read becomes $C=5$ event channels along the coordinate axis:
\begin{itemize}
  \item m6A-marked bases (accessible adenines, the raw Fiber-seq signal);
  \item methylated CpGs (a second, independent chemical mark at C-G letter pairs);
  \item methyltransferase-sensitive patches (MSPs), the accessible stretches;
  \item nucleosomes;
  \item FIRE elements, MSPs classified as regulatory elements.
\end{itemize}
The last three are computed from the m6A pattern. Each channel takes $+1$ (event present), $-1$ (measurable and
absent) or $0$ (not measurable: outside the read, or at a letter where the event cannot occur). The $N$ reads
overlapping a window (we use up to $N=100$) form a tensor of instances $X\in\{-1,0,1\}^{C\times L\times N}$, with a
coverage mask $M\in\{0,1\}^{L\times N}$. The reads are an unordered set. Reads also carry a haplotype tag in phased
samples; the model never sees it, and we use it only in evaluation (Section~\ref{sec:res-xi}).

\textbf{Bulk tracks: the bag labels.} We want $K=8$ other properties on every read: accessibility measured by
ATAC-seq, binding of the protein CTCF, and six histone marks (H3K4me3, H3K27ac, H3K4me1, H3K36me3, H3K27me3,
H3K9me3), which indicate active promoters, enhancers, transcribed gene bodies and two kinds of silenced chromatin.
No technology measures these on many single molecules at once. They are measured in \emph{bulk}: an assay pools
millions of cells of one cell type and reports, per position, how enriched the property is relative to a control
(\emph{fold change}, FC). For cell type $c$ this gives a track $y_{c,k}\in\mathbb{R}_{\ge0}^{L}$ per assay $k$.
ENCODE \citep{ENCODE20} provides these tracks for the cell lines we use, which are also profiled by Fiber-seq. All
cell types share one genome, so every difference between their tracks comes from cell-specific factors, which the
fibers measure and the reference does not.

\textbf{Observation model.} We think of each molecule $f$ as carrying an unobserved per-position value
$u_{f,k}(x)\ge0$, its contribution to assay $k$ at position $x$, in FC units. A bulk track is then the population
average of these values over the molecules covering $x$:
\begin{equation}
y_{c,k}(x)\;=\;\mathbb{E}_{f\sim P_c}\big[\,u_{f,k}(x)\;\big|\;f \text{ covers } x\,\big],
\label{eq:obs}
\end{equation}
where $P_c$ is the population of molecules of cell type $c$. This rests on three assumptions:
\begin{itemize}
  \item[(A1)] a molecule's accessibility, nucleosome and methylation pattern, together with its sequence, carry
    information about its other chromatin properties;
  \item[(A2)] a fold-change track is linear in the fraction of molecules that carry the property, up to a scale
    per experiment;
  \item[(A3)] the Fiber-seq and ENCODE samples of a cell line describe the same population.
\end{itemize}
The fibers of a window are a small sample of $P_c$, and the label averages over the whole population, so the bag
label is not an exact function of the bag.

\textbf{Task.} We learn a function $g_\theta$ that maps one read and the reference to $K$ per-position values,
$z_f=g_\theta(X_{\cdot,\cdot,f},S)\in\mathbb{R}^{K\times L}$, in $\operatorname{asinh}$ FC units, such that
$z_{f,k}(x)\approx\operatorname{asinh}u_{f,k}(x)$. Reads are processed independently: no information passes
between the instances of a bag. The stacked output $Z\in\mathbb{R}^{K\times L\times N}$ is the deliverable, a
per-molecule imputation of every assay. Training sees only the 2-D labels $y_{c,k}$ ($K\times L$ per window); the
model's per-molecule outputs are combined into a bag prediction that can be compared with them
(Section~\ref{sec:aggregation}). The model must generalize to positions not seen in training (a held-out
chromosome) and to cell types whose tracks were never used.

\textbf{The null model.} A model can fit every bag label without using any molecule-specific information: it
predicts the bulk average of the read's cell type on every read,
\begin{equation}
H_0:\quad z_{f,k}(x)=\operatorname{asinh}y_{c,k}(x)\quad\text{for every read } f \text{ of cell type } c.
\label{eq:null}
\end{equation}
Under $H_0$ the bag prediction is perfect, but the per-molecule outputs carry no information beyond the cell type.
Bulk accuracy is therefore necessary but not sufficient, and our hypothesis $H$, that the per-molecule outputs carry
information specific to each molecule, has to be tested with readouts whose value under $H_0$ is known
(Section~\ref{sec:evaluation}).

\begin{table}[t]
\caption{The data for one window of $L=5{,}000$ positions ($x$ the genomic position, $N\le100$ reads, $C=5$
input channels, $K=8$ assays).}
\label{tab:data}
\small
\begin{tabular}{@{}llll@{}}
\toprule
Object & Shape & Values & Role \\
\midrule
Reference $S$ & $4\times L$ & one-hot & input, all cell types \\
Reads $X$ & $C\times L\times N$ & $\{-1,0,1\}$ & input, per molecule \\
Coverage $M$ & $L\times N$ & $\{0,1\}$ & reads covering $x$ \\
Output $Z$ & $K\times L\times N$ & asinh FC & per-molecule output \\
Bag pred.\ $\hat Y$ & $K\times L$ & asinh FC & aggregate of $Z$ \\
Track $Y$ & $K\times L$ & FC $\ge0$ & bag label (bulk) \\
\bottomrule
\end{tabular}
\end{table}

\section{Method}
\label{sec:method}

SMIF has three parts: a network applied to each read on its own (Section~\ref{sec:model}), an aggregation that
turns its per-read outputs into a bulk prediction (Section~\ref{sec:aggregation}), and a loss on that bulk
prediction (Section~\ref{sec:training}).
% TODO: overview figure (reads + reference -> per-read network -> Z -> FC-space average -> loss), referenced here
% and in the architecture paragraph.

\subsection{Per-molecule model}
\label{sec:model}

\textbf{Inputs.} SMIF reads four of the five channels of Section~\ref{sec:problem}: m6A, CpG methylation, MSPs and
nucleosomes. We leave out FIRE. FIRE elements are called by a classifier trained on GM12878 DNase-seq and CTCF
ChIP-seq peaks, so the channel brings bulk-assay knowledge into the input. Leaving it out costs nothing in bulk
accuracy and gives a stronger and more specific single-molecule CTCF read (Section~\ref{sec:res-ctcf},
Appendix~\ref{app:ablations}). Each
read is scored over its \emph{span}: its alignment, trimmed to its first and last nucleosome call. Inputs are taken
over the whole alignment, but the model's outputs, the coverage count and the loss use only the span ($M$ in
Table~\ref{tab:data}).

\textbf{Architecture.} The network maps one read and the window's reference to $K$ tracks of length $L$:
\begin{itemize}
  \item \emph{Encoders.} Two convolutional blocks with the same design, one for the read's channels and one for
    the one-hot reference, each map their input to 8 features per position. Each block has two convolutions of
    width 15 with group normalization, GELU activations and a residual connection. The reference is encoded once
    per window and shared by all its reads. The two feature maps are concatenated.
  \item \emph{U-Net.} Three encoder levels (24, 48 and 96 features, each a convolutional block followed by
    max-pooling by 2) reduce the 5,000 positions to 625 tokens. Three decoder levels (transposed convolutions and
    blocks, with skip connections from the encoder) bring them back to full resolution.
  \item \emph{Transformer bottleneck.} Between the two halves, two transformer encoder layers (96 dimensions,
    4 heads, sinusoidal positional encoding) attend over the 625 tokens, so every output position can use the
    whole 5\,kb of the read.
  \item \emph{Output.} A $1\times1$ convolution gives $K=8$ values per position, and a softplus makes them
    non-negative: $z_f\in\mathbb{R}_{\ge0}^{K\times L}$, in $\operatorname{asinh}$ FC units. Outputs outside the
    read's span are set to 0.
\end{itemize}
The reads of a window share the network's weights and the reference features, but nothing else: there is no
attention or pooling across reads. A read's output therefore depends only on the read and the reference, and the
model is equivariant to the order of the reads. Bulk information enters only through the loss.

\subsection{Fold-change-consistent aggregation}
\label{sec:aggregation}

The label exists only for the bag, so the per-read outputs are combined into a bulk prediction. We read
$z_{f,k}(x)$ as the $\operatorname{asinh}$ of the fold change that a population of molecules like $f$ would show.
Under assumption (A2) a bulk track is an average over molecules in FC units (Eq.~\ref{eq:obs}), so we average in FC
space and map the result back:
\begin{equation}
\hat y_k(x)=\operatorname{asinh}\Big(\frac{1}{|F_x|}\sum_{f\in F_x}\sinh z_{f,k}(x)\Big),
\qquad F_x=\{f: M_{x,f}=1\},
\label{eq:agg}
\end{equation}
where $F_x$ is the set of reads whose span covers $x$. Three properties follow:
\begin{itemize}
  \item if all reads agree, $\hat y_k(x)=z_{f,k}(x)$, so a read's output is on the same scale as the bulk track;
  \item a read with a high value moves the bulk prediction as much as the same signal would move a real bulk
    assay, rather than being damped by averaging in a compressed scale;
  \item a bag that pools the reads of two cell types in proportions $q_A,q_B$ has, under (A2), the exact label
    $\operatorname{asinh}(q_A\sinh t_A+q_B\sinh t_B)$, where $t=\operatorname{asinh}y$. The mixture test
    (Section~\ref{sec:eval-mixture}) relies on this.
\end{itemize}
Working in $\operatorname{asinh}$ units compresses the range of the tracks (fold changes reach about 40 at strong
CTCF peaks) while staying linear near 0, where most positions lie. For numerical stability, $z$ is clipped at 15
before the $\sinh$; real targets stay below about 8.

\subsection{Loss and training}
\label{sec:training}

\textbf{Loss.} The loss is the squared error between the bulk prediction and the track, in $\operatorname{asinh}$
units, averaged over the scored entries:
\begin{equation}
\mathcal{L}=\frac{\sum_{k,x} w_{k}(x)\,\big(\hat y_k(x)-\operatorname{asinh}y_k(x)\big)^2}{\sum_{k,x} w_k(x)}.
\label{eq:loss}
\end{equation}
The weight $w_k(x)=a_k\cdot\mathbf{1}\big[|F_x|\ge5\big]$ keeps only positions covered by at least 5 reads, where
the average of the bag is not dominated by a single read, and only the assays that the cell type has in ENCODE
($a_k\in\{0,1\}$).

\textbf{Training data.} Each training example is one window of one development cell type: its reads (up to 100)
and its 8 tracks. Windows are 5\,kb on the training chromosome, chr20. Half are centred on ENCODE candidate
cis-regulatory elements (cCREs), with a random shift of up to 200\,bp, and half are placed uniformly at random, so
that the model sees both regulatory elements and the background. Validation uses a fixed set of 200 windows on
chr22. The test chromosome, chr21, is used neither for training nor for model selection.

\textbf{Optimization.} We train with AdamW (learning rate $10^{-4}$, batch size 8) for 15 epochs of 5,000 examples,
and reduce the learning rate when the validation loss has not improved for 2 epochs. We use mixed precision and
gradient checkpointing; one training run takes about 2.2 hours on a 20\,GB slice of an NVIDIA H100 GPU. We use the
model after the last epoch. To measure the variability between trainings, every model in this paper is trained
twice, with different random seeds (Section~\ref{sec:res-seeds}).

\textbf{Input variants.} We also train models that differ from SMIF only in their input channels; everything else,
including the data, the architecture and the training budget, is the same. One of them sees only CpG methylation,
so it cannot see accessibility on a molecule; it serves as a null in the single-molecule tests
(Section~\ref{sec:eval-biology}). The others, and what they show about which molecular features carry each assay,
are in Appendix~\ref{app:ablations}.

\section{Evaluation framework}
\label{sec:evaluation}

Our hypothesis $H$ is that the per-molecule outputs carry information specific to each molecule; the null
$H_0$ (Eq.~\ref{eq:null}) predicts the bulk of the read's cell type on every read. We test $H$ against $H_0$ at
three levels. Bulk recovery (Sections~\ref{sec:eval-bulk}--\ref{sec:eval-ceiling}) is necessary but cannot
separate the two. Mixture-only supervision (Section~\ref{sec:eval-mixture}) separates two populations of molecules
without ever showing the model either population's own label. Tests from known biology
(Section~\ref{sec:eval-biology}) check variation \emph{within} one population, between molecules that share a
cell type and a bulk track. Every readout has a known value under $H_0$.

\subsection{Bulk recovery against a cross-cell baseline}
\label{sec:eval-bulk}

We score bulk predictions on the test chromosome, chr21: all 6,843 tiles of 5\,kb, at the positions covered by at
least 5 reads. For each cell type and assay, the predicted track $p$ (the bag prediction $\hat y_k$ of
Eq.~\ref{eq:agg}) is compared with the ENCODE track $t=\operatorname{asinh}y_k$, pooling all scored positions. We
report the mean squared error, the Pearson correlation $r_{pt}$, and the \emph{cell-type correlation} described
below.

Much of every track is shared across cell types, so a strong baseline needs no information about the evaluated
cell type: the \emph{cross-cell average} $b$, at each position the fold-change mean of the development cell types'
tracks, excluding the evaluated cell type's own line, mapped to $\operatorname{asinh}$ units
\citep{Schreiber20Pitfall}. To isolate what is specific to the cell type, we use the partial correlation of $p$ and
$t$ given $b$:
\begin{equation}
r(p,t\mid b)=\frac{r_{pt}-r_{bt}\,r_{pb}}{\sqrt{(1-r_{bt}^2)(1-r_{pb}^2)}},
\label{eq:ctr}
\end{equation}
the correlation of $p$ and $t$ after both are regressed on $b$. It is 1 when the model predicts exactly where and
how the cell type departs from the average, 0 for the baseline itself or any rescaling of it, and negative when the
predicted departures point the wrong way.

We group the cell types into the four \emph{development} lines used in training and three \emph{held-out} lines
never used in training (Section~\ref{sec:data}).

\subsection{Replicate ceiling}
\label{sec:eval-ceiling}

ENCODE tracks are themselves noisy, so no prediction can reach $r=1$. To put all assays on one scale, we estimate
how high any prediction could score. For each of the 54 cell type $\times$ assay targets, we take two biological
replicates of the same ENCODE experiment and score one against the other with the same metrics, at the same
positions. With replicate agreement $\rho$, the pooled target averages $k$ replicates ($k=2$ for 42 targets, 3 for
12) and has Spearman--Brown reliability $k\rho/(1+(k-1)\rho)$. The best correlation any prediction can reach against
it, the \emph{ceiling}, is the square root of this reliability. The same construction on the residuals after the
baseline gives a ceiling for the cell-type correlation. We report each model's \emph{fraction of the ceiling}, model
$r$ divided by ceiling $r$ per target, averaged over targets.

\subsection{Mixture-only supervision}
\label{sec:eval-mixture}

The main test of $H$ trains a model only on bags that mix the reads of two development cell types $A$ and $B$.
The fraction of $A$ is drawn uniformly from $[0.2,0.8]$ per bag, and the label is the fold-change mixture of the two
tracks, $\operatorname{asinh}(q_A\sinh t_A+q_B\sinh t_B)$, exact under (A2) (Section~\ref{sec:aggregation}). Neither
cell type's own track is ever used. We then run each cell type's reads alone. Under $H_0$ the model has only learned
the mixture, so $p_A\approx p_B$; under $H$ each read keeps its cell type's signal. On \emph{differential bins},
1\,kb bins of chr21 where $|t_A-t_B|\ge0.5$, we measure:
\begin{itemize}
  \item \emph{fiber-origin AUC}: in each bin, each read is scored by its mean output over the bin, and the AUC is
    the probability that a random read of the cell type with the higher track scores above a random read of the
    other, averaged over bins (0.5 under $H_0$). It asks whether a \emph{single molecule's} output reveals where
    it came from; even a perfect model stays below 1, because not every molecule of a cell type carries its mark;
  \item \emph{sign accuracy}: the fraction of bins in which the predicted difference $p_A-p_B$ has the sign of
    $t_A-t_B$ (0.5 under $H_0$);
  \item \emph{differential correlation}: the per-position correlation of $p_A-p_B$ with $t_A-t_B$ (0 under $H_0$).
\end{itemize}
The reference for each pair is the same architecture trained on the two cell types' own tracks. As a control for
library effects, we score two Fiber-seq libraries of the same cell line at the same bins; they share one true
track, so any separation between them comes from the libraries, not biology.

\subsection{Single-molecule tests from known biology}
\label{sec:eval-biology}

Mixture supervision separates populations. To test variation \emph{within} a population, we use two properties
known to differ between the molecules of one cell type. Each test follows three rules: the property is never a
label in training; the expected direction of every readout was written down before the run; and the test contains
its own null, a set of molecules or loci where the expected difference is zero. Both tests run SMIF on every read
over a set of target loci and summarize each read's output and measured inputs over the target.

\textbf{CTCF footprints.} CTCF is a protein that binds a specific 15\,bp DNA motif. On a single Fiber-seq read, a
bound CTCF protects its motif from the methyltransferase, leaving a gap in the m6A pattern inside an accessible
patch: a \emph{footprint} \citep{Stergachis20}. The model sees only bulk CTCF ChIP-seq, so it was never told which
molecules are bound.
\begin{itemize}
  \item \emph{Sites.} CTCF motif hits on chr21 (JASPAR MA0139.2, $p\le10^{-5}$). \emph{Bound} sites are the best hit
    inside each ENCODE CTCF peak of a development line (330 sites); in a given cell line, a site counts as bound if
    that line has a peak there. \emph{Unbound} sites are hits more than 500\,bp from every development line's peak,
    sampled to match the bound sites' motif scores in 1-bit bins. We drop unbound sites where the bulk CTCF or ATAC
    track has no reads at all over $\pm1$\,kb in some development line (74 of 272): there, ``no peak'' means
    ``not measured'', not ``unbound''.
  \item \emph{Fiber states}, from the measured input only: \emph{nucleosome} if at least half of the motif is
    covered by a nucleosome call; \emph{footprint} if at least half of it lies in an MSP and its m6A rate is at most
    half that of the 100\,bp flanks; \emph{open} if it lies in an MSP without that protection.
  \item \emph{Hypotheses.} (H1) At bound sites, footprint reads get more predicted CTCF than open reads at the same
    site (per-site difference of means, Wilcoxon signed-rank test over sites with at least 3 reads in each state).
    (H2) The difference is larger for CTCF than for ATAC, since both states are accessible. (H3) Nucleosome-covered
    reads get the least CTCF. (H4, data only) Footprints are more frequent among accessible reads at bound than at
    unbound sites. (H5) Unbound sites get less CTCF than bound sites of the same motif score.
  \item \emph{Null.} The CpG-only variant (Section~\ref{sec:training}; no m6A, MSP or nucleosome channel) cannot see the
    footprint; it must give zero for every per-read difference.
\end{itemize}

\textbf{X inactivation.} In female cells, one of the two X chromosomes is silenced early in development. At a
silenced gene, the inactive copy (Xi) has its CpG-island promoter methylated, its chromatin closed, the active marks
(ATAC, H3K4me3, H3K27ac) missing, and the repressive marks H3K27me3 and H3K9me3 enriched. Escape genes and genes in
the pseudoautosomal regions (PAR) are active on both copies, and XIST is active only on the Xi. GM12878 is female
and trio-phased, so each read is assigned to one parental haplotype. The model was trained on bulk tracks, in which
the two copies are averaged, on chr20 only, and never sees the haplotype.
\begin{itemize}
  \item \emph{Which copy is the Xi.} The haplotype on which the XIST promoter is unmethylated in the reads
    (measured, not predicted): CpG methylation 0.09 against 0.36.
  \item \emph{Score.} Per gene, $d=$ mean output over Xa reads $-$ mean over Xi reads, over the promoter
    (TSS $\pm500$\,bp), for genes with at least 5 reads on each haplotype.
  \item \emph{Gene groups.} \emph{Silenced} genes (chrX genes outside the PAR and not on a list of known escape
    genes) with a CpG-island promoter, 425 genes: the test group. X inactivation silences a promoter by methylating
    its CpG island; silenced genes without one show no measured difference between the copies, so they test
    nothing. \emph{Escape and PAR} genes (33) and chr21 genes (211, where the two haplotypes are an arbitrary pair)
    are nulls. XIST is the reversed case.
  \item \emph{Hypotheses.} (H1) At silenced promoters, $d>0$ for ATAC, H3K4me3 and H3K27ac and $d<0$ for H3K27me3
    and H3K9me3. (H2) Escape and PAR genes have smaller $|d|$. (H3) chr21 genes have $d\approx0$. (H4) XIST is the
    reverse of H1.
  \item \emph{Assays not tested.} H3K4me1 (depleted at the TSS itself and enriched on the flanks, so its direction
    in a promoter window is unclear), H3K36me3 (a gene-body mark, low at promoters) and CTCF (bound at specific
    motifs, not at promoters in general) have no expected direction in a promoter window.
\end{itemize}
The reads' methylation and accessibility already differ between the copies at silenced promoters, so this test
does not show the model \emph{discovering} the Xi. It shows whether the model turns those differences into the
right direction for each mark, on a chromosome and haplotypes it never saw.

\section{Results}
\label{sec:results}

% Preliminary results (Nov 2026 runs). Unless stated otherwise, numbers are the mean of the two training seeds.

\subsection{Data}
\label{sec:data}

We use Fiber-seq data from seven cell lines with ENCODE fold-change tracks for the eight assays (H1 and H9 have no
ATAC track, so 54 targets in total). The four development lines are GM12878 (B lymphoblastoid), K562 (leukemia),
HepG2 (liver) and H1 (embryonic stem cells). The three held-out lines are H9 (embryonic stem cells), HCT116 (colon)
and MCF-7 (breast), ordered from near to far from the development set. Training uses chr20, validation chr22, and
every result below is on chr21.
% TODO: library details (enzyme dose, coverage per line), ENCODE accessions -> Appendix B.

\subsection{Bulk recovery}
\label{sec:res-bulk}

\begin{table*}[t]
\caption{Bulk recovery on chr21, mean over assays and cell types of each group (SMIF: mean of two training seeds).
``Ceiling'' columns divide each correlation by its replicate ceiling (Section~\ref{sec:eval-ceiling}); the ceiling
on $r$ is 0.74 in both groups. The baseline is the cross-cell average, whose cell-type correlation is 0 by
construction.}
\label{tab:bulk}
\small
\begin{tabular}{@{}lccccccc@{}}
\toprule
& \multicolumn{2}{c}{Pearson $r$} & \multicolumn{2}{c}{$r$ / ceiling} & Cell-type $r$ & Cell-type $r$ / ceiling &
\multicolumn{1}{c}{MSE} \\
\cmidrule(lr){2-3}\cmidrule(lr){4-5}
Cell types & SMIF & baseline & SMIF & baseline & SMIF & SMIF & SMIF / baseline \\
\midrule
Development (4) & \textbf{0.50} & 0.44 & \textbf{0.66} & 0.59 & 0.37 & 0.54 & \textbf{0.22} / 0.29 \\
Held out (3) & 0.44 & \textbf{0.48} & 0.59 & \textbf{0.65} & 0.28 & 0.42 & \textbf{0.18} / 0.20 \\
\bottomrule
\end{tabular}
\end{table*}

Table~\ref{tab:bulk} summarizes bulk recovery. ENCODE replicates of the same experiment agree only at $r\approx0.40$
on chr21, so the reachable correlation is about 0.74. In the development lines, SMIF reaches two thirds of that,
ahead of the cross-cell baseline on every metric, and its cell-type correlation (0.37) shows that it predicts where
each cell type departs from the average. In held-out cell types it is behind the baseline on Pearson $r$ but ahead
on MSE, and its cell-type correlation stays positive (0.28, 42\% of its ceiling). The held-out ceilings are as high
as the development ones, so the drop is the model's, not noisier targets. It is concentrated in three assays: the
held-out cell-type correlation is 0.35--0.56 for ATAC, H3K4me3, H3K27ac and H3K4me1, but 0.20 for H3K27me3, 0.14 for
CTCF and 0.06 for H3K36me3.
% TODO: per-assay table or figure (fraction of ceiling per assay and group) in the main text or the appendix.

\subsection{Mixture-only supervision}
\label{sec:res-mixture}

\begin{table}[t]
\caption{Mixture-only supervision on two pairs of development cell types: fiber-origin AUC / sign accuracy /
differential correlation, mean over the pair's assays (K562--H1 has no ATAC). Values under $H_0$: 0.5 / 0.5 / 0.
These models use all five input channels.}
\label{tab:mixture}
\footnotesize
\begin{tabular}{@{}lcc@{}}
\toprule
Trained on & GM12878--HepG2 & K562--H1 \\
\midrule
GM12878 + HepG2 mixtures & \textbf{0.637} / 0.761 / 0.329 & 0.599 / 0.590 / 0.238 \\
GM12878, HepG2 own tracks & 0.635 / 0.771 / 0.351 & 0.600 / 0.598 / 0.224 \\
K562 + H1 mixtures & 0.621 / 0.706 / 0.263 & \textbf{0.666} / 0.802 / 0.291 \\
K562, H1 own tracks & 0.617 / 0.719 / 0.270 & 0.665 / 0.828 / 0.335 \\
\bottomrule
\end{tabular}
\end{table}

A model trained only on mixtures separates the two cell types' molecules as well as a model trained on their own
tracks, on both pairs (Table~\ref{tab:mixture}). On GM12878 and HepG2 the fiber-origin AUC is 0.637 against 0.635,
and on K562 and H1 it is 0.666 against 0.665. Sign accuracy and differential correlation are slightly lower (by 0.01--0.04), and
each cell type's own track is recovered as well, within 0.015 in $r$. Under $H_0$ the mixture model would give both
cell types the mixture, and all three readouts would sit at their null values. Two libraries of the same cell line
separate far less than the real pairs: their mean $|\mathrm{AUC}-0.5|$ is 0.04--0.09, against 0.17--0.25. The
separation is therefore not a library effect. On the pair it did not see, a model does equally well with either kind of supervision; what matters
is having seen the cell types.
% TODO: rerun the mixture test with SMIF's inputs (no FIRE); these models use all five channels.

\subsection{CTCF footprints on single molecules}
\label{sec:res-ctcf}

\begin{table}[t]
\caption{CTCF footprints on chr21: per-site median of (mean predicted CTCF on footprint reads $-$ mean on open
reads), GM12878 / K562, for two training runs of SMIF and for the null model, which sees only CpG methylation
(run 1). ``CTCF $-$ ATAC'' subtracts the same difference for predicted ATAC (specificity: both states are
accessible). Value under $H_0$: 0.}
\label{tab:ctcf}
\footnotesize
\begin{tabular}{@{}lccc@{}}
\toprule
& \multicolumn{2}{c}{SMIF} & CpG only \\
\cmidrule(lr){2-3}
Readout & run 1 & run 2 & (null) \\
\midrule
Bound, footprint $-$ open & \textbf{+1.14 / +0.84} & \textbf{+0.95 / +0.68} & 0.00 / 0.00 \\
Bound, CTCF $-$ ATAC & \textbf{+0.67 / +0.54} & \textbf{+0.85 / +0.63} & 0.00 / 0.00 \\
Unbound, footprint $-$ open & $-0.07$ / $-0.06$ & $-0.03$ / $-0.05$ & $-0.01$ / 0.00 \\
\bottomrule
\end{tabular}
\end{table}

The data first: among accessible reads, 41--43\% carry a footprint at bound sites, against 10\% at unbound sites,
the background rate of our simple call (H4). Table~\ref{tab:ctcf} gives the model results.
\begin{itemize}
  \item \emph{SMIF reads the footprint (H1).} At bound sites, footprint reads get more predicted CTCF than open reads
    at the same site: +1.14 in GM12878 and +0.84 in K562 (94\% and 83\% of sites; Wilcoxon $p<10^{-22}$), and
    +0.95 / +0.68 in the second training. Nucleosome-covered reads get the least (H3).
  \item \emph{The response is specific to CTCF (H2).} Predicted ATAC rises less on footprint reads, so the CTCF
    difference exceeds the ATAC one by +0.54 to +0.85.
  \item \emph{The nulls hold.} At unbound motifs the footprint difference is zero or slightly negative. The CpG-only
    model, which cannot see the footprint, gives 0.00. At the site level, bound sites get more predicted CTCF than
    unbound ones of the same motif score in every score bin (H5; 1.63 against 0.49 in GM12878).
  \item \emph{Bulk accuracy does not predict this read.} A model trained identically but with FIRE as an extra input
    has the same bulk accuracy (development $r$ 0.49 against 0.50), yet its footprint response is ten times weaker
    in one of its two trainings (+0.11 / +0.04) and much less specific in both (Appendix~\ref{app:ablations}). Only a
    single-molecule test can tell the two models apart.
\end{itemize}

On individual reads (Figure~\ref{fig:loci}), SMIF gives footprint reads a block of high CTCF about 200\,bp wide,
centred on the motif and bounded by the accessible patch. Its height on footprint reads comes closer to the bulk
CTCF track at the motif (73--89\% of the bulk level at two of three example sites) than on open reads (37--48\%).
The bulk ChIP peak itself is about twice as wide, because ChIP-seq spreads each binding event over its fragment
length.
% TODO: Figure fig:loci (locus figures, plots/26_nov_20/loci/: one bound CTCF site with the bulk track overlaid,
% the unbound control, and the X-inactivation promoters PIGA and XIST); the figure files must be copied next to
% the .tex (plots/ is gitignored).

\subsection{X inactivation per haplotype}
\label{sec:res-xi}

\begin{table}[t]
\caption{X inactivation in GM12878: fraction of the 425 silenced CpG-island promoters where Xa $-$ Xi has the
expected sign, SMIF training run 1 / run 2 (0.5 under $H_0$).}
\label{tab:xi}
\small
\begin{tabular}{@{}llc@{}}
\toprule
Assay & Expected & SMIF \\
\midrule
ATAC & Xa higher & 0.90 / 0.90 \\
H3K4me3 & Xa higher & 0.87 / 0.86 \\
H3K27ac & Xa higher & 0.90 / 0.86 \\
H3K9me3 & Xa lower & 0.71 / 0.65 \\
H3K27me3 & Xa lower & 0.46 / 0.17 \\
\bottomrule
\end{tabular}
\end{table}

The reads confirm the biology first: at silenced CpG-island promoters, the Xi copy is more methylated (median Xa
$-$ Xi in CpG methylation $-0.23$) and less accessible (MSP fraction $+0.08$), while escape, PAR and chr21 genes
show no difference. Table~\ref{tab:xi} gives the model results.
\begin{itemize}
  \item \emph{Active marks (H1): pass.} SMIF puts ATAC, H3K4me3 and H3K27ac higher on the active copy at 86--90\%
    of silenced promoters in both trainings (median $d=+0.3$ to $+0.6$, $p<10^{-55}$).
  \item \emph{Nulls (H2, H3): pass.} Escape and PAR genes, and chr21 genes, show no difference ($|$median $d|\le0.03$,
    43--64\% in the ``expected'' direction); silenced genes differ 2--6 times more than escape and PAR genes, for
    every mark.
  \item \emph{XIST (H4): pass.} The active marks are reversed in both trainings (H3K4me3: $d=-1.31$ and $-0.87$),
    and H3K9me3 is higher on the Xa, the silent copy of XIST.
  \item \emph{H3K9me3: weak.} SMIF gets the direction right in both trainings (71\% and 65\%), but the differences
    are small (median $d=-0.02$).
  \item \emph{H3K27me3: fails.} SMIF predicts \emph{more} H3K27me3 on the active copy (46\% and 17\% in the expected
    direction), and so does every input variant (Appendix~\ref{app:ablations}). The predicted H3K27me3 follows
    whichever copy's promoter is open and unmethylated: the Xa at silenced genes, but the Xi at XIST.
\end{itemize}
On autosomes, promoter H3K27me3 marks Polycomb targets, which are unmethylated, nucleosome-depleted CpG islands:
the same features that mark the active copy here. On the Xi, H3K27me3 covers methylated, closed islands, a
combination absent from the training chromosome. Part of the failure may also be the test's: at a silenced gene
that is not expressed in GM12878 at all, an unmethylated island on the Xa can be a genuine Polycomb target. The
clearest failure in our locus figures, PTCHD1, is a neuronal gene with almost no active marks on either copy.
Restricting the test to genes expressed in GM12878, and scoring broad Xi domains rather than promoters, will
separate the two.

\subsection{Variability between trainings}
\label{sec:res-seeds}

We trained every model twice. Bulk results barely move: one training changes a cell type $\times$ assay correlation
by a median of 0.010, and a group mean by 0.003--0.013, so every bulk conclusion above holds in both trainings.
Single-molecule readouts keep their direction in both trainings, but their size varies more: SMIF's CTCF footprint
response by 1.2--1.3$\times$, and that of the input variants by up to 10$\times$ (Appendix~\ref{app:ablations}).
Molecule-level claims therefore need several trainings, and we report both throughout.

\section{Discussion}
\label{sec:discussion}
% TODO: what the model reads from a molecule; limitations (bulk ChIP as the only label, seed dependence, H3K27me3
% on the Xi, promoter windows, provisional escape list, chr20-only training); next tests.

\section{Conclusion}
\label{sec:conclusion}
% TODO

\section*{Disclosure of AI use}
% TODO: adjust the scope to what the final version actually used AI for.
We used Claude (Claude Opus 5.5, Anthropic), through the Claude Code tool, as an assistant throughout this project.
It helped write and debug code for the data pipeline, model training and analyses; ran and summarized analyses on
our compute cluster under our direction; and helped draft and edit text, including this manuscript. All
experimental designs, hypotheses and conclusions were decided by the authors, who checked the results and the text
and take full responsibility for the content of this paper.

\bibliographystyle{ACM-Reference-Format}
\bibliography{proposal}

\appendix
\section{Biology background}
\label{app:bio}
% TODO: from proposal.tex appendix A.

\section{Data and training details}
\label{app:data}
% TODO: sample list, ENCODE accessions, configs (data24, model16, train04), compute.

\section{Input ablations}
\label{app:ablations}
% TODO: to be written. Draft content moved here from the main text (results section and CTCF / X-inactivation
% comparisons across models); rewrite as a self-contained appendix.

\begin{table}[t]
\caption{Input ablations: Pearson $r$ in the development cell types on chr21, mean of two training runs. Each model
differs from SMIF only in its input channels. One training moves these values by about 0.003--0.013.}
\label{tab:ablations}
\small
\begin{tabular}{@{}lccccc@{}}
\toprule
Assay & SMIF & All inputs & m6A + CpG & CpG only & No CpG \\
\midrule
ATAC & 0.77 & 0.76 & 0.77 & 0.62 & 0.76 \\
H3K4me3 & 0.52 & 0.53 & 0.53 & 0.48 & 0.49 \\
H3K27ac & 0.51 & 0.49 & 0.51 & 0.40 & 0.46 \\
H3K4me1 & 0.56 & 0.53 & 0.57 & 0.40 & 0.52 \\
CTCF & 0.40 & 0.39 & 0.39 & 0.31 & 0.39 \\
H3K36me3 & 0.48 & 0.48 & 0.49 & 0.47 & 0.24 \\
H3K27me3 & 0.45 & 0.44 & 0.45 & 0.41 & 0.36 \\
H3K9me3 & 0.39 & 0.39 & 0.39 & 0.30 & 0.39 \\
\midrule
Mean & 0.50 & 0.49 & 0.50 & 0.42 & 0.44 \\
\bottomrule
\end{tabular}
\end{table}

We trained four models that differ from SMIF only in their inputs (Table~\ref{tab:ablations}): \emph{all inputs}
(SMIF plus FIRE), \emph{m6A + CpG} (no fibertools calls: no MSP, nucleosome or FIRE), \emph{CpG only} (no
accessibility information at all; the null of the CTCF test) and \emph{no CpG} (all inputs except CpG
methylation).
\begin{itemize}
  \item \emph{The fibertools calls add nothing in bulk.} SMIF and m6A + CpG are level, and adding FIRE costs
    0.02--0.03 on H3K27ac and H3K4me1 in both runs. The model gets what it needs from the raw m6A and CpG calls.
  \item \emph{m6A carries the accessibility-type assays.} Without it, ATAC, H3K4me1, H3K27ac, CTCF and H3K9me3 drop
    by 0.09--0.16 in development lines, and more in held-out lines (ATAC $-0.21$, H3K4me1 $-0.22$, CTCF $-0.15$).
  \item \emph{CpG methylation carries H3K36me3 and H3K27me3.} Without it H3K36me3 halves ($-0.24$), consistent with
    gene-body methylation following H3K36me3, and H3K27me3 drops by 0.09. H3K4me3 uses both inputs.
\end{itemize}

\begin{table*}[t]
\caption{CTCF footprints at bound sites for every input variant: per-site median of (mean predicted CTCF on
footprint reads $-$ mean on open reads), GM12878 / K562, for both training runs.}
\label{tab:ctcf-ablations}
\small
\begin{tabular}{@{}lccccc@{}}
\toprule
& \multicolumn{2}{c}{Footprint $-$ open, CTCF} & \multicolumn{2}{c}{CTCF $-$ ATAC} & Unbound sites,
footprint $-$ open \\
\cmidrule(lr){2-3}\cmidrule(lr){4-5}
Model & run 1 & run 2 & run 1 & run 2 & run 1 \,\textbar\, run 2 \\
\midrule
SMIF & \textbf{+1.14 / +0.84} & \textbf{+0.95 / +0.68} & \textbf{+0.67 / +0.54} & \textbf{+0.85 / +0.63} &
$-0.07$ / $-0.06$ \,\textbar\, $-0.03$ / $-0.05$ \\
All inputs & +0.11 / +0.04 & +0.73 / +0.44 & +0.02 / +0.05 & +0.23 / +0.21 & $-0.02$ / $-0.01$ \,\textbar\, $-0.06$ /
$-0.08$ \\
m6A + CpG & +0.26 / +0.20 & +0.56 / +0.58 & +0.20 / +0.14 & +0.55 / +0.50 & $-0.05$ / $-0.03$ \,\textbar\, $-0.01$ /
$+0.02$ \\
No CpG & +0.18 / +0.10 & +0.75 / +0.46 & $-0.12$ / $-0.02$ & +0.33 / +0.31 & $-0.08$ / $-0.08$ \,\textbar\,
$-0.09$ / $-0.10$ \\
CpG only (null) & 0.00 / 0.00 & 0.00 / 0.01 & 0.00 / 0.00 & 0.00 / 0.00 & $-0.01$ / 0.00 \,\textbar\, $-0.01$ /
$-0.01$ \\
\bottomrule
\end{tabular}
\end{table*}

\textbf{CTCF footprints across inputs} (Table~\ref{tab:ctcf-ablations}). Every model that sees m6A reads the
footprint, but the size varies 2--10$\times$ between trainings for the three variants other than SMIF, and their
specificity is low in at least one training. A likely reason for the difference between SMIF and the all-inputs
model: with FIRE in the input, a channel derived from CTCF ChIP peaks, the model can take ``this is a CTCF site''
from that channel instead of reading the molecule. Without it, it has to read the m6A pattern, and it learns, from
bulk supervision alone, that a protected motif on an accessible molecule means CTCF.

\textbf{X inactivation across inputs.} Every variant puts the active marks higher on the active copy (82--91\% of
silenced promoters in both trainings), including the CpG-only model, so either input alone suffices. H3K9me3 is
stable only without CpG input (86\% in both trainings) and varies between trainings with it (21--76\%). H3K27me3
fails in every variant (17--35\%).

\end{document}

\endinput
